\subsection{Evaluation and decision support}
\label{subsec:evaluator}

The last step is the evaluation of the simulated plans. At this point no new
traffic is generated and no additional SUMO simulations are run. The task is
to take the repeated outputs from the Monte Carlo experiment and compare the
candidate detours in a way that is fair, reproducible and useful for the
final discussion. The key point is to keep the paired structure of the
experiment: for each Monte Carlo seed, the evaluator has one baseline result
and one result for each candidate deviation plan. Each plan is compared
against the baseline from the same run, not against an unrelated average.

In practice, the evaluation is organised around three questions:
\begin{itemize}
  \item How much each plan changes the network-level performance compared with the baseline.
  \item Whether the plan remains feasible, in particular whether it avoids reducing the number of completed trips.
  \item Among the feasible plans, which alternatives still look reasonable after Pareto screening and a sensitivity check on the decision weights.
\end{itemize}

\subsubsection{Choice of evaluator metrics}
\label{sssec:evaluator_metric_choices}

The evaluator does not try to prove that one detour is universally best. It
separates the comparison into quantities that answer different planning
questions:
\begin{itemize}
  \item \textbf{Total time-loss delay} measures the network-level congestion
        burden created or removed by a plan. It is the primary efficiency
        metric because it captures both queues and slow movement relative to
        the vehicle's ideal route traversal.
  \item \textbf{Mean travel time} measures the average trip burden for
        completed vehicles. It complements time-loss delay because a plan can
        reduce congestion on the chosen route while still making paths longer,
        or conversely add local congestion while shortening many trips.
  \item \textbf{Completed and unfinished trips} are used as a feasibility
        diagnostic. A plan must not appear good merely because more vehicles
        fail to finish within the simulation window and are therefore absent
        from completed-trip averages.
  \item \textbf{Fraction of impacted users delayed} focuses on the vehicles
        whose route actually changed. Network averages can hide a plan that
        performs well overall but imposes most of the cost on the directly
        affected users.
\end{itemize}

These metrics are also practical: they are available for every Monte Carlo run
from SUMO's \texttt{tripinfo.xml} output and from the route comparison step, so
the same definitions can be applied consistently to every candidate plan. The
statistical tests reported by the evaluator are diagnostic checks on the
paired effects, they are not the ranking rule. The final decision layer uses
effect sizes, feasibility, Pareto dominance and weight sensitivity.

\subsubsection{Comparable run set}
\label{sssec:evaluator_comparable_runs}

The first operation is to identify the runs that can actually be compared.
The evaluator scans the Monte Carlo output directory for
\texttt{run\_\#/results.json} files and keeps a run only when it contains the
baseline and all candidate plans listed in \texttt{summary.json}. Runs with
missing plan outputs are left out of the paired analysis, but they are still
recorded in the metadata so that the exclusion is visible.

This step avoids comparing one plan over more replications than another. After
the filtering, every paired statistic is computed on the same run identifiers
for every plan. If no complete run is available, the evaluator stops instead
of producing a ranking based on incomplete evidence.

\subsubsection{Paired network-level statistics}
\label{sssec:evaluator_paired_stats}

For every plan $\mathcal{D}_k$ and every retained Monte Carlo run $r$, the
basic quantity is the difference between the plan and the baseline from the
same run. For the primary network-level metric, this gives
\begin{equation}
  \Delta D_k^{(r)} = D_k^{(r)} - D_0^{(r)},
  \label{eq:evaluator_delta_delay}
\end{equation}
where $D_0^{(r)}$ is the baseline time-loss delay from run $r$ and
$D_k^{(r)}$ is the corresponding value after applying plan $\mathcal{D}_k$.
For mean travel time, the same construction is used:
\begin{equation}
  \Delta \bar{\tau}_k^{(r)}
  =
  \bar{\tau}_k^{(r)} - \bar{\tau}_0^{(r)}.
  \label{eq:evaluator_delta_travel_time}
\end{equation}

The sign convention is kept simple: a positive paired difference means that
the plan worsened the metric in that run and a negative value means that it
improved it. The full vector of paired differences is kept, not only the
average, because the consistency of an effect matters as much as its mean.

For a generic paired-difference vector
$x_1,\ldots,x_n$, the evaluator reports
\begin{equation}
  \bar{x},\quad s,\quad \min(x),\quad \max(x),\quad \tilde{x},
  \label{eq:evaluator_descriptive_stats}
\end{equation}
where $\bar{x}$ is the sample mean, $s$ the sample standard deviation and
$\tilde{x}$ the median. The mean estimates the expected effect of the plan
over the seed distribution, the standard deviation describes run-to-run
variability and the median is kept as a robust centre when a few runs are
extreme.

The confidence interval answers a different question: how precisely the
Monte Carlo experiment has estimated the mean paired effect. Because the true
standard deviation of the paired effects is unknown and must be estimated
from the $n$ replications, the normal critical value is replaced by the
Student-\textit{t} critical value. The Student-\textit{t} distribution has
heavier tails than the normal distribution, so it gives wider intervals when
only a small number of runs is available. The evaluator uses the standard
95\% level, corresponding to a two-sided 5\% error rate, as a conventional
compromise between intervals that are too narrow to be reliable and intervals
that are too wide to be useful. The interval is computed as
\begin{equation}
  \bar{x}
  \pm
  t_{0.975,\,n-1}\frac{s}{\sqrt{n}},
  \label{eq:evaluator_t_interval}
\end{equation}
when at least two runs are available. In repeated experiments satisfying the
usual independence assumptions, this procedure would contain the true mean
effect in about 95\% of such intervals, it does not mean that there is a
95\% probability that this particular realized interval contains the truth.

Two diagnostic hypothesis tests are reported for each paired metric. Both ask
whether the observed paired effects are compatible with a zero-effect
baseline. The paired one-sample \textit{t}-test checks the null hypothesis
\begin{equation}
  H_0: \mathbb{E}[\Delta M_k] = 0,
\end{equation}
where $\Delta M_k$ is the paired effect of plan $k$ on the metric being
tested. It is most appropriate when the paired differences are approximately
normally distributed. A Wilcoxon signed-rank test is also reported as a
non-parametric companion. It ranks the absolute non-zero paired differences,
attaches the original signs and checks whether positive and negative ranks are
balanced around zero. This makes the check less dependent on the normality
assumption, which is useful for traffic simulation outputs with small samples
or occasional extreme runs.

\subsubsection{Feasibility gate}
\label{sssec:evaluator_feasibility}

Before plans are ranked, they pass through a completed-trip feasibility check.
A candidate is considered feasible only if it completes at least as many trips
as the baseline in every retained run:
\begin{equation}
  C_k^{(r)} \geq C_0^{(r)}
  \qquad
  \forall r \in \{1,\ldots,n\},
  \label{eq:evaluator_feasibility}
\end{equation}
where $C_k^{(r)}$ is the number of completed trips under plan
$\mathcal{D}_k$. The reason for this rule is straightforward: a detour should
not look better simply because more vehicles failed to arrive before the end
of the simulation window. Plans that fail the gate are still shown in the
report, but they are not allowed into the Pareto screening or weighted score.

\subsubsection{Impacted-vehicle metrics}
\label{sssec:evaluator_impacted_users}

Network-level metrics are necessary, but they can hide what happens to the
drivers who are directly affected by the closure. For this reason, the
evaluation also considers an impacted vehicle set. When route files are
available, a vehicle is counted as impacted for a plan if its edge sequence
differs between the baseline route file and the plan-specific route file:
\begin{equation}
  \mathcal{I}_k =
  \{\,v : R_{k,v} \neq R_{0,v}\,\}.
  \label{eq:evaluator_impacted_set}
\end{equation}
If route-difference information is not available, the evaluator can fall back
to the vehicles whose baseline route crosses the closed edge. This is less
precise, but it still keeps the analysis focused on users directly connected
to the disruption.

For each retained run, the impacted-user analysis compares only vehicles that
are present and completed in both the baseline and plan \texttt{tripinfo.xml}
files. It reports the total additional travel time for impacted users,
\begin{equation}
  \Delta T_{k,\mathcal{I}}^{(r)}
  =
  \sum_{v \in \mathcal{I}_k}
  \left(\tau_{k,v}^{(r)} - \tau_{0,v}^{(r)}\right),
  \label{eq:evaluator_impacted_delta_time}
\end{equation}
the mean effective speed over their trajectories and the fraction of impacted
users who are delayed:
\begin{equation}
  \phi_k^{(r)}
  =
  \frac{
    |\{\,v \in \mathcal{I}_k :
      \tau_{k,v}^{(r)} > \tau_{0,v}^{(r)}\,\}|
  }{
    |\mathcal{I}_k|
  }.
  \label{eq:evaluator_phi}
\end{equation}
The mean of $\phi_k^{(r)}$ across runs is used as one of the three decision
criteria. Lower values are preferred, because they mean that a smaller share
of directly affected users experiences a travel-time increase.

\subsubsection{Multi-criteria decision analysis}
\label{sssec:evaluator_mcda}

The final ranking is built in three steps. First, infeasible plans are removed
by the completed-trip gate of Section~\ref{sssec:evaluator_feasibility}.
Second, the remaining plans are screened for Pareto dominance using three cost
criteria:
\begin{equation}
  \left(
    \overline{\Delta D_k},
    \overline{\Delta \bar{\tau}_k},
    \overline{\phi_k}
  \right).
  \label{eq:evaluator_criteria_vector}
\end{equation}
A feasible plan $a$ Pareto-dominates a feasible plan $b$ if it is no worse on
all three criteria and strictly better on at least one. Dominated plans are
removed before the weights are applied. This prevents the final ranking from
selecting a plan that is worse than another available plan on all measured
dimensions.

Third, the plans left on the Pareto front are ranked with a normalized
additive score. For each criterion $q$, values are min-max normalized over the
Pareto front:
\begin{equation}
  z_{kq}
  =
  \frac{x_{kq} - \min_j x_{jq}}
       {\max_j x_{jq} - \min_j x_{jq}},
  \label{eq:evaluator_normalisation}
\end{equation}
with $z_{kq}=0$ when all surviving plans have the same value on criterion
$q$. The composite score is then
\begin{equation}
  S_k = \sum_q w_q z_{kq},
  \label{eq:evaluator_additive_score}
\end{equation}
where all criteria are costs and lower scores are better. The reference
configuration gives equal weight to network time-loss delay, mean travel time and the
fraction of impacted users delayed. This should be read as a neutral starting
point, not as a claim that the three dimensions are inherently equally
important.

Because the weights contain judgement, the evaluator also checks how sensitive
the ranking is to them. It samples weight vectors uniformly from the
three-criterion simplex using a Dirichlet$(1,1,1)$ distribution. Here a
weight vector is a triple $(w_1,w_2,w_3)$ with non-negative entries that sum
to one. For each sampled weight vector, the evaluator
recomputes the ranking. The report then shows how often each
Pareto-efficient plan ranks first, its mean rank, score quantiles and
pairwise win probabilities. A plan that wins under the reference weights and
also remains strong under many sampled weights is more convincing than a plan
that only wins under a narrow weighting assumption.